\section{Results}
\label{sec:results}

\subsection{Averaging embeddings names the field a cluster shares}
\label{sec:cluster-labeling}

\begin{figure}[t]
\centering
\includegraphics[width=\linewidth]{pacs-clustering.pdf}
\caption{
    Cluster labeling and mixing with \doctolora.
    \textbf{(a)} Fuzzy token-set overlap between each decoded label and official 28 PACS labels. Each point is a node, with distribution represented by kernel density estimation (KDE).
    \textbf{(b)} Pairwise judge win rates, showing the fraction of PACS labels where the generated label from the row method is judged closer to the official label than that from the column. Color runs blue (row loses), grey (tie), to orange (row wins), row means at right.
    \textbf{(c)} Cluster centroids, decoded per PACS field/division/subdivision. Horizontal axis is vector length, circle area is the number of papers. Decoded labels (orange) and official PACS names (grey, if different) are shown.
    \textbf{(d)} Specificity as a function of vector length: longer vectors yield specific, shorter more general concepts. For each Wikipedia seed, lines show continuity of label, grey bands mark collapse to base-model prior.
    \textbf{(e)} Decoded abstract position (SBERT space) as a function of B weight, A$\to$B, for 50 near pairs. \doctolora and \texttt{ICAE} track the target, in-context stays near midpoint, \texttt{T2L} flat.
    \textbf{(f)} Verbatim copy rate by method, measured as the fraction of output tokens present in sources. \texttt{ICAE} is highest, and \doctolora and in-context prompting are comparable.
}
\label{fig:cluster-labels}
\end{figure}

A widespread practice for intepretating clustering results is either manual interpretation of the members, or running a keyword or LLM pipeline that needs careful tuning before it returns a usable name~\citep{grootendorst2020keybert,grootendorst2022bertopic,grootendorst2023keyllm,murray2025labeling}.
With \doctolora, labeling becomes as simple as averaging the vectors of a cluster and decoding the mean (\figref{cluster-labels}c).

We validate cluster labeling on physics papers from the American Physical Society (APS), using the Physics and Astronomy Classification Scheme (PACS) (Appendix~\ref{app:data}).
PACS is a hierarchical classification system for topics in physics, in which each topic is assigned a numeric code. 
For example, ``03.65'' is assigned to quantum mechanics.
If the last two digits are omitted, it refers to the broader, higher-level topic above it. For example, ``03'' refers to quantum mechanics and quantum field theory, while ``0'' refers to general physics.
Authors assign codes to their papers to organize them, with the first code indicating the paper's main topic.
We group papers based on their first code at each of the three levels of the hierarchy (field, division, subdivision) and calculate the average of the vectors within each group.
We decode these averages using the instruction ``In 2 to 3 words, name the scientific field that all of these documents belong to'' and compare the decoded labels with the official PACS names of the nodes (Appendix~\ref{app:baseline-labels}).

We compare \doctolora to four baselines: \texttt{KeyLLM}~\citep{grootendorst2023keyllm}, which prompts an LLM with the original documents to extract keywords. 
\texttt{vec2text}~\citep{morris2023vec2text}, which iteratively edits text to match a target embedding.
\texttt{ICAE}~\citep{ge2024icae}, which encodes documents into memory slots and decodes them back.
\texttt{T2L}~\citep{charakorn2025texttolora}, which generates LoRA adapters from a frozen \texttt{GTE} encoder~\citep{li2023gte}. 
We also compare against the base model's own hidden states (see Appendix~\ref{app:actpatch} for details). 
All methods receive the same input: unmodified paper titles and abstracts.
%In addition to the architectural difference,  \texttt{T2L} is trained to reconstruct task adapters rather than given text, whereas \doctolora is trained to decode the text directly.

We quantify the quality of the decoded labels by fuzzy token-set overlap, a score from $0$ to $1$ with no LLM in the loop (\figref{cluster-labels}a).
\doctolora leads with $.642 \pm .055$, ahead of \texttt{KeyLLM}, \texttt{vec2text}, \texttt{ICAE}, and \texttt{T2L}.
As an additional validation, we use five LLM judges (\texttt{glm-5.3-flash}, \texttt{gemini-3.8-flash}, \texttt{gpt-5.6-luna}, \texttt{muse-spark-1.3}, and \texttt{mistral-medium-3.1}) to compare every pair of methods on each node (\figref{cluster-labels}b).
LLM judges consistently prefer \doctolora's labels over all baselines.
A control rules out that any cluster draws a plausible field name: on clusters sampled at random across PACS chapters the judges accept no label at all, against $73.3\%$ for the real nodes, and rescaling the control mean to the length of a real node does not change that (Appendix~\ref{app:incoherent}).

%It produces three or four generic strings unrelated to physics across all 28 nodes.

Upon manual inspection, \doctolora decodes cluster means into formal field names such as ``particle physics'' and ``condensed matter physics'' more frequently than all baselines (\figref{cluster-labels}c).
The baselines exhibit four main failure modes.
\texttt{T2L} produces a handful of unrelated, fixed field names (e.g., ``Social psychology'') for all nodes; \texttt{vec2text} generates indecipherable pseudo-words (e.g., ``inferromagnetic spin-diabetic scattering interactions''). \texttt{KeyLLM} returns overly specific keyword lists from few medoid papers. \texttt{ICAE} decodes plausible field names labels are too broad (``Materials science'') or too specific (``Quantum many-body systems''). 
The differing performance of \doctolora and \texttt{T2L} shows that just creating a LoRA adapter is not enough. 
Decodable clustering depends on how the adapter is trained.

\doctolora avoids generating cluster names that are too broad or too narrow because it interprets the length of the mean vector as an indicator of abstraction level. 
When document vectors point in the same direction (i.e., are well aligned), their average keeps much of the original length. 
If they vary widely, averaging pulls the mean closer to zero, shortening the vector. For \doctolora, we observe that the average vector length grows as we move from general categories (top-level fields) to more detailed ones (fine-grained subdivisions; \figref{cluster-labels}c).
This effect appears in all methods, but only \doctolora can turn vector length into the appropriate level of abstraction when decoding names. 
To test this, we embed and deliberately rescale Wikipedia introductions (\figref{cluster-labels}d). 
With high scaling ($\alpha$), we get detailed concepts such as ``Citric acid cycle metabolism.'' 
As we reduce $\alpha$, the decoded label shifts to something more general, like ``Cellular respiration.'' 
For Nikola Tesla, shrinking the vector gives ``AC power.'' For a light-emitting diode, it shifts to ``Semiconductor device structure'' (Appendix~\ref{app:length-dial}). 
Extremely short vectors produce generic names from the base model, as if the LoRA adapter were absent. 
If the vector is too long, the output fails to decode because it drifts off the space of meaningful representations.

\subsection{Midpoints decode to blends of both sources}
\label{sec:fusion}

\begin{table}[t]
    \centering
    \footnotesize
    % three wide decoded-text columns ran ~3pt past \textwidth at the default
    % column padding; trimming \tabcolsep pulls them back inside the margin
    \setlength{\tabcolsep}{4pt}
    \caption{One near pair decoded along the edge at $25$, $50$ and $75\,\%$ weight on paper B, under the prompt quoted in the text.
    Paper A is ``Statistical mechanics of learning from examples'' (PhysRevA 45.6056) and paper B is ``Good quantum error-correcting codes exist'' (PhysRevA 54.1098).
    Each cell is the opening of the decoded abstract, whole sentences cut at $22$ words.
    Where the decoded abstracts actually land between the two papers, averaged over all $50$ near pairs at each of the $13$ grid points, is \figref{cluster-labels}e.
    }
    \label{tab:mixing-decode}
    % Generated by workflow/scripts/fig2_mixing_table.py (snakemake fig2); do not edit by hand.
% Pair pairCSML_00: A = Statistical mechanics of learning from examples (10.1103/PhysRevA.45.6056); B = Good quantum error-correcting codes exist (10.1103/PhysRevA.54.1098).
% Cells: opening of the decoded abstract (whole sentences past 12 words, cut at 22);
\begin{tabular}{@{}l *{3}{>{\raggedright\arraybackslash}p{0.275\linewidth}}@{}}
\toprule
& \multicolumn{3}{c}{Ideal mixing weight on paper B} \\
\cmidrule(lr){2-4}
Method & $25\,\%$ & $50\,\%$ & $75\,\%$ \\
\midrule
\doctolora & This research investigates the learning of real-valued functions in neural networks using a statistical-mechanical framework, focusing on the generalization error of trained\ldots{} & This research investigates the existence and properties of quantum error-correcting codes in the context of quantum computing, focusing on the asymptotic behavior\ldots{} & This research topic investigates the existence and properties of quantum error-correcting codes, which are essential for protecting quantum information from decoherence\ldots{} \\
\texttt{ICAE} & This study investigates the learning of statistical mechanics from feedforward neural networks, focusing on the existence of stochastic, Gibbs-free energy-based models\ldots{} & This research explores the existence of quantum error-correcting codes in the context of quantum mechanics\ldots{} & This research explores the existence of quantum error-correcting codes (QECCs) in the context of quantum computing\ldots{} \\
in-context & This research investigates the statistical-mechanical behavior of quantum neural networks under the influence of both stochastic learning dynamics and quantum decoherence, blending\ldots{} & This research investigates the statistical-mechanical behavior of quantum neural networks under the influence of both stochastic learning dynamics and quantum decoherence, blending\ldots{} & This research investigates the emergence of quantum spin-glass phases in neural network architectures under the influence of quantum decoherence, blending statistical mechanics\ldots{} \\
\bottomrule
\end{tabular}

\end{table}

\doctolora can decode text from any point in its embedding space, including points between documents. 
We test this by interpolating between APS paper pairs with controlled PACS taxonomy distance: 50 ``near'' pairs from the same subtopic, and 50 ``far'' pairs from different subtopics, both sampled randomly.
For each pair, we decode at 13 points along the interpolation between documents using \doctolora, in-context prompting, \texttt{ICAE}, and \texttt{T2L}, all under the prompt \textit{``Write a detailed abstract (four to six sentences) describing this research topic: its problem, methods, and findings.''} (\tabref{mixing-decode}; pair sampling, per-method interpolation, and metrics in Appendix~\ref{app:edge-protocol}).
In-context prompting has an advantage, as it directly receives both endpoint documents and mixing proportions, while \doctolora, \texttt{ICAE}, and \texttt{T2L} only receive source content indirectly via parameters, e.g., LoRA weights or memory slots. 
We paraphrase this instruction three times and decode every edge again under each wording.
The separation between the methods is unchanged.
The slope of the decoded position against the requested weight stays at $1.11$--$1.20$ for \doctolora and $1.19$--$1.24$ for \texttt{ICAE}, against $0.29$--$0.41$ for in-context prompting, and no wording moves a slope by more than $0.04$ (Appendix~\ref{app:prompt-sensitivity}).

We evaluate similarity to endpoints using SBERT-based scores.
We also calculate the verbatim copy rate~\citep{grusky2018newsroom}, which is the fraction of output tokens shared with either source.
While \doctolora and \texttt{ICAE} follow the true mixing weights across the edge, in-context prompting remains near the midpoint, except at the endpoints (\figref{cluster-labels}e).
The copy rate separates \texttt{ICAE} from both alternatives.
\texttt{ICAE} reuses source wording at $0.73$ on either stratum, while \doctolora ($0.56$) and in-context prompting ($0.52$--$0.55$) stay close to each other (\figref{cluster-labels}f).
\texttt{T2L} collapses onto a few input-independent texts.
Its decoded position stays within $0.490$--$0.506$ across all $13$ points, endpoints included, while \doctolora moves from $0.081$ to $0.904$ (\figref{cluster-labels}e).
The endpoints are therefore not decodable, and interpolation is moot.


\subsection{Retrieval with \doctolora embeddings}
\label{sec:similarity}

\doctolora is designed to induce specific behaviors in LLMs, which may or may not contain topic-level information about each document.
If \doctolora adapters do encode topics within their representations, and if we can extract and align this information geometrically, \doctolora has practical value as an embedding for retrieval tasks.
To test the capacity of \doctolora embeddings for retrieval, we apply a simple linear transformation to the \doctolora embedding space, aiming to make topic similarity more directly accessible while preserving decodability, and then evaluate its effectiveness on standard retrieval benchmarks.

We adapt \doctolora embeddings with a lightweight, linear transformation $g_\theta$ (262k parameters for Qwen), which is trained (on $42{,}332$ uniformly-sampled OpenAlex citation edges, less than $0.002\%$ of $2.56$B possible pairs) to rebalance the representation for similarity tasks, while, crucially, without sacrificing decodability.
We compare \doctolora embeddings with and without $g_\theta$ against commonly used text encoders---\texttt{SPECTER2}~\citep{singh2023scirepeval}, \texttt{Instructor}~\citep{su2023instructor}, \texttt{EmbeddingGemma}~\citep{embeddinggemma2025}, \texttt{SBERT}~\citep{reimers2019sbert}, \texttt{GTE}~\citep{li2023gte} on the following tasks (Appendix~\ref{app:data} for data details and Appendix~\ref{app:benchmarks} for the full benchmark protocol).

\begin{table}[t]
\centering
\small
\caption{Similarity benchmarks, absolute scores. Each row is one task on one field or dataset; \emph{raw} is the \doctolora embedding and \emph{adapted} applies $g_\theta$, the linear transform trained on OpenAlex. Text encoders are ordered by mean rank, best first. Cells are the bootstrap mean $\pm$ its standard deviation over $1{,}000$ resamples. The best method in a row is bold and shaded and the second-best underlined. Rows carry different metrics and sit at different levels, so a method is read down a row, never across rows. \texttt{SPECTER} replaces \texttt{SPECTER2} on the disambiguation rows. Full protocol in Appendix~\ref{app:benchmarks}.}
\label{tab:similarity}
% Auto-generated by workflow/plot/tab_similarity_benchmarks.py -- do not edit by hand.
\setlength{\tabcolsep}{4pt}
\renewcommand{\arraystretch}{1.05}
\resizebox{\textwidth}{!}{%
\begin{tabular}{@{}llccccccc@{}}
\toprule
 & & \multicolumn{2}{c}{\doctolora} & \multicolumn{5}{c}{text encoders} \\
\cmidrule(lr){3-4}\cmidrule(l){5-9}
Task & Field / dataset & raw & adapted & \texttt{SBERT} & \texttt{GTE} & \texttt{Emb.Gemma} & \texttt{Instructor} & \texttt{SPECTER2} \\
\midrule
\multirow{3}{*}{Next-paper (AUC)}
 & Economics & .813\std{.001} & \cellcolor{gray!10}\underline{.910}\std{.001} & \cellcolor{gray!25}\textbf{.937}\std{.001} & \cellcolor{gray!10}\underline{.910}\std{.001} & .906\std{.001} & .877\std{.001} & .903\std{.001} \\
 & Psychology & .799\std{.001} & \cellcolor{gray!10}\underline{.923}\std{.001} & \cellcolor{gray!25}\textbf{.937}\std{.001} & .912\std{.001} & .910\std{.001} & .888\std{.001} & .908\std{.001} \\
 & Physics & .880\std{.001} & \cellcolor{gray!10}\underline{.958}\std{.001} & .956\std{.001} & \cellcolor{gray!25}\textbf{.962}\std{.001} & .947\std{.001} & .908\std{.001} & .941\std{.001} \\
\midrule
\multirow{3}{*}{Topic (macro-F1)}
 & Economics & .242\std{.015} & .374\std{.022} & \cellcolor{gray!25}\textbf{.418}\std{.022} & \cellcolor{gray!10}\underline{.398}\std{.023} & .362\std{.021} & .379\std{.021} & .344\std{.021} \\
 & Psychology & .217\std{.010} & .340\std{.020} & \cellcolor{gray!25}\textbf{.366}\std{.023} & .350\std{.019} & \cellcolor{gray!10}\underline{.358}\std{.020} & .348\std{.026} & .319\std{.020} \\
 & Physics & \cellcolor{gray!10}\underline{.590}\std{.012} & \cellcolor{gray!10}\underline{.590}\std{.009} & .573\std{.011} & .587\std{.010} & .588\std{.010} & .574\std{.010} & \cellcolor{gray!25}\textbf{.598}\std{.007} \\
\midrule
\multirow{3}{*}{Collaboration (AUC)}
 & Economics & .596\std{.011} & .629\std{.012} & \cellcolor{gray!25}\textbf{.655}\std{.011} & .627\std{.012} & \cellcolor{gray!10}\underline{.633}\std{.012} & .620\std{.012} & .632\std{.011} \\
 & Psychology & .602\std{.012} & \cellcolor{gray!10}\underline{.687}\std{.010} & \cellcolor{gray!25}\textbf{.691}\std{.011} & .668\std{.011} & .666\std{.011} & .649\std{.011} & .665\std{.011} \\
 & Physics & .792\std{.004} & .795\std{.004} & .832\std{.003} & \cellcolor{gray!10}\underline{.839}\std{.003} & .828\std{.003} & \cellcolor{gray!25}\textbf{.854}\std{.003} & .826\std{.004} \\
\midrule
\multirow{5}{*}{Author disamb.\ (B\textsuperscript{3} F1)}
 & zbMATH & .933\std{.001} & .932\std{.001} & \cellcolor{gray!25}\textbf{.944}\std{.001} & \cellcolor{gray!10}\underline{.939}\std{.001} & .934\std{.002} & .936\std{.001} & .934\std{.001} \\
 & QIAN & .787\std{.005} & .850\std{.004} & \cellcolor{gray!10}\underline{.865}\std{.004} & \cellcolor{gray!25}\textbf{.866}\std{.004} & .832\std{.004} & .832\std{.004} & .832\std{.004} \\
 & ArnetMiner & .678\std{.004} & \cellcolor{gray!25}\textbf{.708}\std{.005} & .698\std{.005} & .698\std{.004} & \cellcolor{gray!10}\underline{.706}\std{.004} & .679\std{.004} & .669\std{.004} \\
 & PubMed & .693\std{.007} & .815\std{.006} & \cellcolor{gray!25}\textbf{.832}\std{.006} & .792\std{.006} & \cellcolor{gray!10}\underline{.822}\std{.006} & .801\std{.006} & .739\std{.006} \\
 & KISTI & .708\std{.002} & .767\std{.002} & \cellcolor{gray!25}\textbf{.793}\std{.002} & \cellcolor{gray!10}\underline{.785}\std{.002} & .771\std{.002} & .754\std{.002} & .746\std{.002} \\
\bottomrule
\end{tabular}
}

\end{table}

With $g_\theta$, \doctolora sits inside the band of the text encoders on all four retrieval tasks. It performs better than \texttt{Instructor} and \texttt{SPECTER2}, is comparable to \texttt{EmbeddingGemma} and \texttt{GTE}, and follows \texttt{SBERT}, which leads on 9 of 14 benchmarks. The average ranks are: \texttt{SBERT} $2.0$, \texttt{GTE} $2.9$, \doctolora{+}$g_\theta$ $3.4$, \texttt{EmbeddingGemma} $3.5$, \texttt{Instructor} $4.7$, \texttt{SPECTER2} $5.0$, and raw \doctolora $6.5$. All text encoders use default settings without task-specific tuning. We do not claim \doctolora outperforms current encoders, but it retains similarity information comparable to standard methods.

\texttt{T2L} cannot decode the original documents for cluster labeling or interpolation (\secref{cluster-labeling}, \secref{fusion}), yet the frozen encoder it reads, \texttt{GTE}, is a competitive retrieval encoder. On the means, \texttt{GTE} is ahead of \texttt{SBERT} only on the three Physics tasks and on QIAN, and behind it on every Economics and Psychology task and on the other three disambiguation sets (\tabref{similarity}). The bootstrap intervals separate the two on next-paper prediction, where \texttt{GTE} leads on Physics and trails on Economics and Psychology, and on zbMATH, PubMed, and KISTI, where \texttt{SBERT} leads. The remaining differences lie within the intervals. \texttt{T2L} is not trained to reconstruct documents, and the similarity that \texttt{GTE} carries does not by itself make the adapter decodable.

%\begin{table}[t]
%\centering
%\small
%\caption{What an embedding costs. \emph{Params} counts the encoder that must run per document, with the \doctolora row giving the frozen base model that the \doctolora hypernetwork runs on top of. \emph{Dim} is the vector entering the index. \emph{Index} is that vector stored in fp16 for the ${\approx}2.2$M papers of Appendix~\ref{app:data}. \emph{Wins} counts first places over the $14$ tests of \tabref{similarity}. Full-rank \doctolora is the tensor that gets decoded and is never indexed (\secref{doc2lora-embedding}).}
%\label{tab:cost}
%\begin{tabular}{@{}l r r r c@{}}
%\toprule
%Encoder & Params & Dim & Index & Wins \\
%\midrule
%\doctolora (Qwen3-4B), pooled & $4.0$B$+$hypernet & $18{,}432$ & $81$~GB & $2$ \\
%\doctolora (Qwen3-4B), full rank & $4.0$B$+$hypernet & $147{,}456$ & $648$~GB & --- \\
%\midrule
%\texttt{SPECTER2} (\texttt{specter2\_base}) & $110$M & $768$ & $3.4$~GB & $1$ \\
%\texttt{SBERT} (\texttt{all-mpnet-base-v2}) & $110$M & $768$ & $3.4$~GB & $10$ \\
%\texttt{Instructor} (\texttt{instructor-large}) & $335$M & $768$ & $3.4$~GB & $1$ \\
%\texttt{EmbeddingGemma} (\texttt{embeddinggemma-300m}) & $300$M & $768$ & $3.4$~GB & $0$ \\
%\bottomrule
%\end{tabular}
%\end{table}
%